%----------------------------------------------------------------------------------------
% LOGARITMISCHE FUNCTIES, VERGELIJKINGEN EN ONGELIJKHEDEN. Alle opgaven zijn eigen werk.
%
% Opbouw: Start, leerstof (theorie, werkwijze of uitleg, voorbeeld, korte oefening), en op een
% nieuwe pagina de Opdrachten (door elkaar, stijgende moeilijkheid), Koppel en Test jezelf.
% Labels van dit hoofdstuk beginnen met h3-.
%----------------------------------------------------------------------------------------

\mychapter{3}{Logaritmische functies}

\lkr{Leerplan B+S: LPD 7-12 (exponentiële en logaritmische functies), het deel over de logaritmische functie, logaritmische vergelijkingen en ongelijkheden. De inverse functie komt uit de cursus Algemene functies. Het grootste struikelblok zijn de bestaansvoorwaarden: laat ze altijd eerst opschrijven, en de oplossingen op het einde ertegen controleren.}

In hoofdstuk 2 was de logaritme een getal. Nu bekijken we $x \mapsto \log_a x$ als functie. Ze is de inverse van de exponentiële functie: ze draait $a^x$ terug. Daarmee los je ook vergelijkingen en ongelijkheden op waarin de onbekende in een logaritme staat.

\begin{oefeningen}[Start]
\oef{h3-start-1}

\oef{h3-verken}
\end{oefeningen}

%========================================================================================
\section{De logaritmische functie}

\subsection{Definitie, domein en bereik}

In hoofdstuk 2 hoorde bij elk positief getal $x$ één logaritme $\log_a x$. Elk positief getal krijgt zo één beeld: dat is een functie.

\begin{theorie}
Voor $a > 0$ met $a \neq 1$ is $\formule{f(x) = \log_a x}$ de \textbf{\textcolor{accent}{logaritmische functie}} met grondtal $a$.
\end{theorie}

Welke getallen mag je invullen, en welke beelden komen eruit? Dat volgt rechtstreeks uit de definitie van de logaritme.

\begin{theorie}
\begin{stelling}{Domein en bereik}
Voor $f(x) = \log_a x$ is \qquad $\formule{\dom f = \left]0, +\infty\right[}$ \qquad en \qquad $\formule{\ber f = \R}$
\end{stelling}
\begin{bewijs}
De logaritme $\log_a x$ bestaat precies als $x > 0$: dat is het domein. Elk reëel getal $y$ is een beeld, want $a^y > 0$ ligt in het domein en
\bewijsregel{\log_a a^y = y .}
\end{bewijs}
\end{theorie}

Een logaritmische functie neemt dus elke waarde aan, maar je mag er alleen positieve getallen in invullen. Bij een samengestelde functie zoals $\ln(4 - x^2)$ moet daarom alles onder de logaritme strikt positief zijn.

\voorbeeld{}{Bepaal het domein van $f(x) = \ln(4 - x^2)$.\label{vb:h3-domein}}{%
\uitleg{Van een logaritme moet het argument strikt positief zijn.}
\begin{stappenbreed}
4 - x^2 > 0 \iff x^2 < 4 & Bestaansvoorwaarde. \\
\dom f = \left]-2, 2\right[ & Tussen de nulwaarden van de bergparabool.
\end{stappenbreed}}

\begin{oefeningen}
\oef{h3-k-domein}
\end{oefeningen}

\subsection{De inverse van de exponentiële functie}

Volgens de definitie is $y = a^x \iff x = \log_a y$. De functie $\log_a x$ draait dus terug wat $a^x$ doet. Wat betekent dat voor de grafieken?

\begin{theorie}
\begin{stelling}{Inverse functie}
De functie $\log_a x$ is de \textbf{\textcolor{accent}{inverse functie}} van $a^x$. Hun grafieken zijn elkaars spiegelbeeld om de rechte $\formule{y = x}$.
\end{stelling}
\begin{bewijs}
Een punt $(p, q)$ ligt op de grafiek van $a^x$ als $q = a^p$. Volgens de definitie van de logaritme is dat zo precies als $p = \log_a q$, dus als het punt $(q, p)$ op de grafiek van $\log_a x$ ligt. De punten $(p, q)$ en $(q, p)$ zijn elkaars spiegelbeeld om $y = x$.
\end{bewijs}
\end{theorie}

\naast{%
Domein en bereik wisselen van plaats: $2^x$ heeft alleen positieve beelden, dus $\log_2 x$ bestaat alleen voor positieve $x$. Het vaste punt $(0, 1)$ van $2^x$ wordt $(1, 0)$, en $(1, 2)$ wordt $(2, 1)$. De horizontale asymptoot van $2^x$, de $x$-as, wordt na het spiegelen een verticale asymptoot: de $y$-as.

Een logaritmische functie stijgt heel traag: $\log_2 x$ bereikt $10$ pas bij $x = 1024$. Dicht bij $0$ daalt ze onbeperkt: $\log_2 0{,}001 \approx -10$.}{%
\begin{tikzpicture}[x=.5cm,y=.5cm,font=\scriptsize]
  \rooster{-3}{1}{5}{-3}{1}{5}
  \draw[grijs,dashed] (-3,-3) -- (5.2,5.2) node[above,text=black] {$y = x$};
  \kromme[dashed]{-3}{2.32}{2^(\x)}
  \kromme{.13}{5}{ln(\x)/ln(2)}
  \vol{(1,0)} \vol{(2,1)}
  \node[right,text=accent] at (4.2,1.6) {$\log_2 x$};
  \node[left,text=accent] at (1.9,4.6) {$2^x$};
\end{tikzpicture}}

\needspace{14\baselineskip}
\subsection{Stijgen en dalen}

Ook het stijgen of dalen neemt de logaritmische functie over van de exponentiële functie met hetzelfde grondtal.

\begin{theorie}
\begin{stelling}{Stijgen en dalen}
Als $a > 1$, dan is $\log_a x$ \textbf{stijgend}: \qquad $\formule{0 < x < y \Rightarrow \log_a x < \log_a y}$

Als $0 < a < 1$, dan is $\log_a x$ \textbf{dalend}: \qquad $\formule{0 < x < y \Rightarrow \log_a x > \log_a y}$
\end{stelling}
\begin{bewijs}
Neem $a > 1$ en stel dat $\log_a x \geqslant \log_a y$. Omdat $a^t$ stijgt (of bij gelijke exponenten gelijk is), is dan
\bewijsregel{x = a^{\log_a x} \geqslant a^{\log_a y} = y ,}
in strijd met $x < y$. Dus is $\log_a x < \log_a y$. Neem nu $0 < a < 1$ en stel dat $\log_a x \leqslant \log_a y$. Omdat $a^t$ daalt, is dan $x = a^{\log_a x} \geqslant a^{\log_a y} = y$, weer in strijd met $x < y$. Dus is $\log_a x > \log_a y$.
\end{bewijs}
\end{theorie}

\figuren{\begin{tikzpicture}[x=.5cm,y=.5cm,font=\scriptsize]
  \rooster{-1}{1}{6}{-3}{1}{3}
  \kromme{.13}{6}{ln(\x)/ln(2)}
  \vol{(1,0)} \vol{(2,1)}
  \node[above left,text=accent] at (6,2.6) {$\log_2 x$};
  \node[below,text=black] at (2.5,-3.3) {$a = 2 > 1$: stijgend};
\end{tikzpicture}\qquad\qquad
\begin{tikzpicture}[x=.5cm,y=.5cm,font=\scriptsize]
  \rooster{-1}{1}{6}{-3}{1}{3}
  \kromme{.06}{6}{-ln(\x)/ln(3)}
  \vol{(1,0)} \vol{(.333,1)}
  \node[below left,text=accent] at (6,-1.7) {$\log_{\frac13} x$};
  \node[below,text=black] at (2.5,-3.3) {$a = \frac13 < 1$: dalend};
\end{tikzpicture}}

Daaruit volgt meteen het teken van een logaritme. Is $a > 1$, dan is $\log_a x$ negatief voor $0 < x < 1$ en positief voor $x > 1$, want $\log_a 1 = 0$. Bij $0 < a < 1$ is het net omgekeerd.

\voorbeeld{}{Is $\log_{0{,}4} 3$ positief of negatief? Antwoord zonder rekentoestel.\label{vb:h3-teken}}{%
\uitleg{Vergelijk met $\log_{0{,}4} 1 = 0$, en kijk of de functie stijgt of daalt.}
\begin{stappenbreed}
0 < 0{,}4 < 1 & Het grondtal is kleiner dan 1: $\log_{0{,}4} x$ is dalend. \\
3 > 1 \ \Longrightarrow\ \log_{0{,}4} 3 < \log_{0{,}4} 1 = 0 & Dalend: een groter getal geeft een kleinere logaritme.
\end{stappenbreed}
$\log_{0{,}4} 3$ is negatief. Controle: $\log_{0{,}4} 3 = \frac{\ln 3}{\ln 0{,}4} \approx -1{,}199$.}

\begin{oefeningen}
\oef{h3-k-teken}
\end{oefeningen}

\needspace{14\baselineskip}
\subsection{De grafiek}

Uit wat voorafgaat, kennen we nu de belangrijkste kenmerken van de grafiek. We vatten ze samen.

\begin{theorie}
\begin{stelling}{Grafiek van een logaritmische functie}
De grafiek van $\log_a x$ gaat door $\formule{(1, 0)}$ en $\formule{(a, 1)}$, en heeft de $y$-as als \textbf{\textcolor{accent}{verticale asymptoot}}.
\end{stelling}
\end{theorie}

De twee punten volgen uit $\log_a 1 = 0$ en $\log_a a = 1$. De asymptoot is het spiegelbeeld van de horizontale asymptoot van $a^x$; dat de grafiek de $y$-as echt nadert, bewijzen we niet, want daarvoor heb je limieten nodig. Met die kenmerken en de transformaties van grafieken schets je snel een verwante functie.

\voorbeeld{}{Schets de grafiek van $g(x) = \log_2(x - 3) + 1$. Geef het domein, de asymptoot en de nulwaarde.\label{vb:h3-grafiek}}{%
\naast{%
\uitleg{De grafiek van $\log_2 x$ schuift 3 naar rechts en 1 naar boven; de asymptoot schuift mee naar rechts.}
\begin{stappenvrij}{4.6cm}{3cm}
(1, 0) \mapsto (4, 1) & Het vaste punt. \\
\dom g = \left]3, +\infty\right[ & $x - 3 > 0$. \\
\text{asymptoot } x = 3 & \\
\log_2(x - 3) = -1 & Nulwaarde. \\
x - 3 = \tfrac12 \iff x = \tfrac72 & Definitie.
\end{stappenvrij}}{%
\begin{tikzpicture}[x=.5cm,y=.5cm,font=\scriptsize]
  \rooster{-1}{1}{7}{-3}{1}{3}
  \draw[grijs,dashed] (3,-3) -- (3,3.3);
  \kromme{3.07}{7}{ln(\x-3)/ln(2)+1}
  \vol{(4,1)} \vol{(3.5,0)}
\end{tikzpicture}}}

\begin{oefeningen}
\oef{h3-k-schets}

\oef{h3-grafiek}
\end{oefeningen}

%========================================================================================
\needspace{14\baselineskip}
\section{Logaritmische vergelijkingen}

Een logaritmische vergelijking heeft de onbekende onder een logaritme. Je lost ze op met de definitie van de logaritme, of met de volgende eigenschap.

\begin{theorie}
\begin{stelling}{Gelijke logaritmen}
Voor $A, B > 0$ is \qquad $\formule{\log_a A = \log_a B \iff A = B}$
\end{stelling}
\begin{bewijs}
Is $A = B$, dan zijn ook de logaritmen gelijk. Omgekeerd: is $\log_a A = \log_a B$, verhef dan $a$ tot beide leden:
\bewijsregel{A = a^{\log_a A} = a^{\log_a B} = B .}
\end{bewijs}
\end{theorie}

Twee gelijke logaritmen met hetzelfde grondtal hebben dus gelijke argumenten. Maar let op de voorwaarde $A, B > 0$: een oplossing van $A = B$ waarvoor $A$ of $B$ niet positief is, is geen oplossing van de logaritmische vergelijking. Daarom schrijf je altijd eerst de \textbf{bestaansvoorwaarden} op.

\voorbeeld{}{Los op: $\log_3(2x - 1) = 2$.\label{vb:h3-vergelijking}}{%
\uitleg{Schrijf eerst de bestaansvoorwaarde op. Gebruik dan de definitie: de logaritme is de exponent.}
\begin{stappenbreed}
2x - 1 > 0 \iff x > \tfrac12 & Bestaansvoorwaarde. \\
2x - 1 = 3^2 = 9 & Definitie van de logaritme. \\
x = 5 & Voldoet aan $x > \frac12$.
\end{stappenbreed}}

\begin{oefeningen}
\oef{h3-k-def}
\end{oefeningen}

Staat er aan beide kanten een logaritme met hetzelfde grondtal, dan stel je de argumenten gelijk.

\voorbeeld{}{Los op: $\log_2(x^2 - 3) = \log_2(2x)$.\label{vb:h3-gelijk}}{%
\uitleg{Bestaansvoorwaarden eerst, dan de eigenschap van gelijke logaritmen. Controleer op het einde elke oplossing.}
\begin{stappenbreed}
x^2 - 3 > 0 \text{ en } 2x > 0 \iff x > \sqrt3 & Bestaansvoorwaarden. \\
x^2 - 3 = 2x \iff x^2 - 2x - 3 = 0 & Gelijke logaritmen. \\
(x - 3)(x + 1) = 0 \iff x = 3 \text{ of } x = -1 & Ontbinden. \\
x = 3 & $-1$ voldoet niet: $\log_2(-2)$ bestaat niet.
\end{stappenbreed}}

\begin{oefeningen}
\oef{h3-vergelijking}

\needspace{5\baselineskip}
\oef{h3-k-gelijk}
\end{oefeningen}

Staan er meerdere logaritmen in hetzelfde lid, breng ze dan eerst samen met de rekenregels.

\voorbeeld{}{Los op: $\log_2(x + 1) + \log_2(x - 1) = 3$.\label{vb:h3-twee}}{%
\uitleg{Breng het linkerlid samen tot één logaritme met de productregel. Dat mag alleen waar beide logaritmen bestaan: daarom eerst de bestaansvoorwaarden.}
\begin{stappenbreed}
x + 1 > 0 \text{ en } x - 1 > 0 \iff x > 1 & Bestaansvoorwaarden. \\
\log_2\bigl((x + 1)(x - 1)\bigr) = 3 & Productregel. \\
x^2 - 1 = 8 \iff x = 3 \text{ of } x = -3 & Definitie. \\
x = 3 & $-3$ voldoet niet aan $x > 1$.
\end{stappenbreed}}

Waar komt de valse oplossing $-3$ vandaan? Voor $x = -3$ is $(x + 1)(x - 1) = (-2) \cdot (-4) = 8$, dus $\log_2\bigl((x + 1)(x - 1)\bigr) = 3$ klopt. Maar $\log_2(x + 1) = \log_2(-2)$ bestaat niet. De productregel vergroot het domein: het product van twee negatieve getallen is positief. Daarom controleer je de oplossingen altijd tegen de bestaansvoorwaarden van de \textit{oorspronkelijke} vergelijking.

\begin{werkwijze}
Schrijf de \textbf{bestaansvoorwaarden} op: elk argument van een logaritme moet strikt positief zijn. Breng de vergelijking met de rekenregels in de vorm $\log_a A = c$ (dan is $A = a^c$) of $\log_a A = \log_a B$ (dan is $A = B$). Los op en houd alleen de oplossingen die aan de bestaansvoorwaarden voldoen.
\end{werkwijze}

\begin{oefeningen}
\oef{h3-k-twee}
\end{oefeningen}

%========================================================================================
\needspace{20\baselineskip}
\section{Ongelijkheden}

Bij een vergelijking stel je exponenten of argumenten gelijk. Bij een ongelijkheid moet je bovendien weten of het teken blijft of omkeert, en dat hangt af van het grondtal.

\subsection{Ongelijkheden met machten}

Een exponentiële functie met een grondtal groter dan 1 stijgt: bij een grotere exponent hoort een grotere macht. Dat geldt ook omgekeerd, en daarmee los je een ongelijkheid op. Bij een grondtal tussen 0 en 1 is alles net andersom.

\begin{theorie}
\begin{stelling}{Ongelijkheden met machten}
Als $a > 1$, dan is \qquad $\formule{a^p < a^q \iff p < q}$

Als $0 < a < 1$, dan is \qquad $\formule{a^p < a^q \iff p > q}$
\end{stelling}
\begin{bewijs}
Neem $a > 1$. Is $p < q$, dan is $a^p < a^q$, want $a^x$ stijgt. Omgekeerd: stel $a^p < a^q$. Was $p \geqslant q$, dan was $a^p \geqslant a^q$ (gelijk als $p = q$, groter als $p > q$), in strijd met $a^p < a^q$. Dus is $p < q$. Voor $0 < a < 1$ verloopt het bewijs net zo, met een dalende functie: daar keren de ongelijkheden tussen de exponenten om.
\end{bewijs}
\end{theorie}

Zijn beide leden een macht van hetzelfde grondtal, dan vergelijk je dus de exponenten: bij een grondtal groter dan 1 blijft het ongelijkheidsteken, bij een grondtal tussen 0 en 1 keert het om. Voor $\leqslant$ geldt hetzelfde, want bij gelijkheid zijn ook de exponenten gelijk.

\voorbeeld{}{Los op: $\Bigl(\dfrac13\Bigr)^x \geqslant 9$.\label{vb:h3-machtong}}{%
\uitleg{Schrijf 9 als macht van $\frac13$. Het grondtal ligt tussen 0 en 1, dus het teken keert om.}
\begin{stappenbreed}
\Bigl(\frac13\Bigr)^x \geqslant \Bigl(\frac13\Bigr)^{-2} & $\bigl(\frac13\bigr)^{-2} = 3^2 = 9$. \\
x \leqslant -2 & Grondtal kleiner dan 1: het teken keert om.
\end{stappenbreed}
Controle: voor $x = -3$ is $\bigl(\frac13\bigr)^{-3} = 27 \geqslant 9$, voor $x = 0$ is $\bigl(\frac13\bigr)^0 = 1 < 9$.}

\begin{oefeningen}
\oef{h3-k-machtong}
\end{oefeningen}

Is het rechterlid geen mooie macht, neem dan van beide leden een logaritme. Omdat $\ln x$ stijgt, blijft het teken daarbij staan. Maar let op: als je daarna deelt door een negatief getal, zoals $\ln 0{,}8$, keert het teken om.

\voorbeeld{}{Een medicijn wordt per uur met $20$\,\% afgebroken. Vanaf wanneer is er minder dan $10$\,\% van over?\label{vb:h3-ongelijkheid}}{%
\uitleg{Na $t$ uur is er nog een fractie $0{,}8^t$ over. Neem $\ln$ van beide leden; $\ln$ stijgt, dus het teken blijft. Daarna deel je door $\ln 0{,}8 < 0$: het teken keert om.}
\begin{stappenbreed}
0{,}8^t < 0{,}1 & \\
t \ln 0{,}8 < \ln 0{,}1 & $\ln$ van beide leden. \\
t > \frac{\ln 0{,}1}{\ln 0{,}8} \approx 10{,}3 & Delen door een negatief getal.
\end{stappenbreed}
Na ongeveer $10{,}3$ uur is er minder dan $10$\,\% over.}

\begin{oefeningen}
\oef{h3-k-expong}
\end{oefeningen}

\needspace{16\baselineskip}
\subsection{Ongelijkheden met logaritmen}

Ook een logaritmische functie stijgt of daalt. Staat er aan beide kanten een logaritme met hetzelfde grondtal, dan vergelijk je dus de argumenten, en het grondtal beslist of het teken blijft.

\begin{theorie}
\begin{stelling}{Ongelijkheden met logaritmen}
Voor $x, y > 0$:

Als $a > 1$, dan is \qquad $\formule{\log_a x < \log_a y \iff x < y}$

Als $0 < a < 1$, dan is \qquad $\formule{\log_a x < \log_a y \iff x > y}$
\end{stelling}
\begin{bewijs}
Neem $a > 1$. Is $x < y$, dan is $\log_a x < \log_a y$, want $\log_a x$ stijgt. Omgekeerd: stel $\log_a x < \log_a y$. Was $x \geqslant y$, dan was $\log_a x \geqslant \log_a y$, in strijd met wat gegeven is. Dus is $x < y$. Voor $0 < a < 1$ verloopt het bewijs net zo, met een dalende functie.
\end{bewijs}
\end{theorie}

Hier komen de bestaansvoorwaarden er nog bij: de stelling geldt alleen voor positieve $x$ en $y$. Een oplossing van een logaritmische ongelijkheid moet dus aan twee voorwaarden tegelijk voldoen.

\voorbeeld{}{Los op: $\log_3(x - 2) < 2$.\label{vb:h3-logongelijkheid}}{%
\uitleg{Bestaansvoorwaarde eerst. Schrijf dan $2$ als logaritme met grondtal 3; omdat $3 > 1$ blijft het teken.}
\begin{stappenbreed}
x - 2 > 0 \iff x > 2 & Bestaansvoorwaarde. \\
\log_3(x - 2) < \log_3 9 \iff x - 2 < 9 & $2 = \log_3 9$, grondtal $> 1$. \\
x \in \left]2, 11\right[ & Beide voorwaarden samen.
\end{stappenbreed}}

\begin{oefeningen}
\oef{h3-ongelijkheid}
\end{oefeningen}

Met een grondtal tussen 0 en 1 keert het teken om. Vergeet ook dan de bestaansvoorwaarde niet: ze geeft de andere grens van de oplossing.

\voorbeeld{}{Los op: $\log_{\frac12}(x + 3) > -2$.\label{vb:h3-logongklein}}{%
\uitleg{Schrijf $-2$ als logaritme met grondtal $\frac12$. Het grondtal is kleiner dan 1, dus het teken keert om.}
\begin{stappenbreed}
x + 3 > 0 \iff x > -3 & Bestaansvoorwaarde. \\
\log_{\frac12}(x + 3) > \log_{\frac12} 4 & $-2 = \log_{\frac12} 4$, want $\bigl(\frac12\bigr)^{-2} = 4$. \\
x + 3 < 4 \iff x < 1 & Grondtal kleiner dan 1: het teken keert om. \\
x \in \left]-3, 1\right[ & Beide voorwaarden samen.
\end{stappenbreed}
Controle: voor $x = -2$ is $\log_{\frac12} 1 = 0 > -2$, voor $x = 5$ is $\log_{\frac12} 8 = -3 < -2$.}

\begin{werkwijze}
Schrijf bij een logaritme eerst de bestaansvoorwaarden op. Schrijf beide leden als macht of als logaritme met hetzelfde grondtal, of neem van beide leden $\ln$. Vergelijk dan exponenten of argumenten: bij een grondtal groter dan 1 blijft het teken, bij een grondtal tussen 0 en 1 keert het om. Ook delen door een negatief getal keert het teken om. Combineer op het einde met de bestaansvoorwaarden.
\end{werkwijze}

\begin{oefeningen}
\oef{h3-k-logong}
\end{oefeningen}

%========================================================================================
\opdrachtenkop

\begin{oefeningen}
\needspace{13\baselineskip}
\oef{h3-o1}

\oef{h3-o11}

\needspace{14\baselineskip}
\oef{h3-o2}

\needspace{12\baselineskip}
\oef{h3-o12}

\oef{h3-o13}

\oef{h3-o3}

\needspace{10\baselineskip}
\oef{h3-o14}

\needspace{18\baselineskip}
\oef{h3-o15}

\oef{h3-o16}

\oef{h3-o4}

\oef{h3-o19}

\oef{h3-o7}

\needspace{14\baselineskip}
\oef{h3-o6}

\oef{h3-o17}

\oef{h3-o20}

\oef{h3-o5}

\oef{h3-o22}

\oef{h3-o8}

\oef{h3-o18}

\oef{h3-o21}

\oef{h3-o23}

\oef{h3-o10}

\oef{h3-o24}

\oef{h3-o26}

\oef{h3-o9}

\oef{h3-o25}
\end{oefeningen}

%----------------------------------------------------------------------------------------
% KOPPEL: verbanden met andere leerstof
%----------------------------------------------------------------------------------------
\needspace{20\baselineskip}
\opdrachtsectie{Koppel}

\lkr{Algemene functies, Rijen en Integralen: kijk na wat al gezien is. Het pakket legt de volgorde niet vast.}

\begin{oefeningen}
\oef{h3-koppel-1}

\oef{h3-koppel-2}

\oef{h3-koppel-3}
\end{oefeningen}

%----------------------------------------------------------------------------------------
\opdrachtsectie{Test jezelf}

\begin{oefeningen}
\oef{h3-t1}

\oef{h3-t2}

\oef{h3-t3}

\oef{h3-t4}

\oef{h3-t5}

\oef{h3-t6}

\oef{h3-t7}
\end{oefeningen}

\herhaalschema{%
  \fv{oef:h3-t1}{vb:h3-grafiek}{oef:h3-grafiek}%
  \fv{oef:h3-t2}{vb:h3-domein}{oef:h3-k-domein}%
  \fv{oef:h3-t3}{vb:h3-vergelijking}{oef:h3-k-def}%
  \fv{oef:h3-t4}{vb:h3-twee}{oef:h3-k-twee}%
  \fv{oef:h3-t5}{vb:h3-ongelijkheid}{oef:h3-k-expong}%
  \fv{oef:h3-t6}{vb:h3-logongklein}{oef:h3-k-logong}%
  \fv{oef:h3-t7}{vb:h3-grafiek}{oef:h3-o3}}

\ifzelfstudie
\section*{Maak je eigen samenvatting}
Sluit dit boek en schrijf in je schrift: domein, bereik, vaste punten en asymptoot van $\log_a x$ voor $a > 1$ en $0 < a < 1$, het verband met $a^x$, de werkwijze voor logaritmische vergelijkingen met de bestaansvoorwaarden, en wanneer het teken van een ongelijkheid omkeert. Vergelijk daarna met de kaders.
\fi
